\section{Format of the Exam}

\subsection{Logistical Information}

The exam takes place in a lockdown browser on the \href{https://ucldata.atlassian.net/wiki/spaces/ELearningStudentSupport/pages/63340834/WISEflow+AssessmentUCL+guidance+for+students}{WiseFlow} platform.
This means that internet access will be limited to the reference materials that we have allowed to be accessed from the lockdown browser.
The particular sites may change between years, but during the exam you will typically have access to:

\begin{itemize}
  \item A JupyterLite server that, in particular, contains the codebase that you will be working on during the exam.
  \item Any \hyperref[sec:pre-reading-material]{pre-reading material} that was made available before the exam.
  \item A PDF version of the exam paper for those who prefer on-screen reading.
  \item The Python standard library documentation, plus the documentation websites of any external libraries that the code you will be working on uses.
\end{itemize}

You will enter your answers via the WiseFlow interface.
For questions that require written answers, you can simply type your answer into the appropriate answer boxes.
For questions that require you to write code, you can use the JupyterLite server to write and test your code, which provides basic Python syntax checking and a means of running the tests.
Once you are happy with your code, you will need to copy-and-paste the code you wrote into the appropriate answer box on WiseFlow.
Note that you can create ``code blocks'' in the WiseFlow answer box, before pasting your code so that any formatting is retained.

You are required to attend the exam \textbf{in-person} - attending remotely is not an option.

You do not need to bring your own laptop / computer, you will be provided with a machine that has the WiseFlow software installed when you arrive at the exam.

\subsection{Tasks and Available Marks}

The maximum mark for the exam is 50.
You will have 2 hours to complete the exam.

There will be 3 tasks on the exam paper, each worth 25 marks.
You must complete 2 tasks out of the 3 - which tasks you choose to complete is up to you to decide.

\textbf{IMPORTANT}: If you provide answers for more than 3 tasks, you will be required to indicate \textbf{which two} tasks you would like to have contribute to your total.
For example; if I complete all 25 marks available in task 1, then answer 10 marks of questions in task 2, and 15 marks of questions in task 3, the mark I will obtain for the exam is:

\begin{itemize}
  \item $35 ( = 25 + 10)$ if I indicate I would like tasks 1 and 2 to contribute.
  \item $40 ( = 25 + 15)$ if I indicate I would like tasks 1 and 3 to contribute.
  \item $25 ( = 10 + 15)$ if I indicate I would like tasks 2 and 3 to contribute.
\end{itemize}

For this reason, we \textbf{strongly recommend} you read through each of the tasks at the start of the exam, before deciding which one(s) to attempt.

\subsection{Pre-Reading Material} \label{sec:pre-reading-material}

Approximately one month before the date of the exam, pre-reading material will be released.
This will contain some contextual information about the setting of the problem you'll be working on during the exam.
You are not required to do any additional reading to support the pre-reading - any information you need about specific algorithms / problems / terminology (etc) pertaining to the problem is provided.
Note however, that the pre-reading material (and indeed the exam!) assumes that you are familiar with the fundamentals of the course.

The pre-reading material will also be available to you during the exam, for reference.
Some of the tasks may require you to refer to the pre-reading material, or make use of information provided within it to complete them.
